\begin{theorem}[Poisson--Cauchy \(f\)-散度喷流的全局奇次消隐律]\label{thm:conclusion-poisson-cauchy-fdiv-jet-odd-order-vanishing}
设 \(R\ge 2\)。假定 \(\tilde\nu\) 具有足够多的中心矩，使得
\[
\delta_t(x):=\frac{\tilde h_t(x)}{\tilde g_t(x)}-1
=
\sum_{k=2}^{2R}\frac{\mu_k}{t^k}U_k(y)+o(t^{-2R}),
\qquad
y:=\frac{x-\bar\gamma}{t},
\]
在 \(G(y)\,\dd y\) 的加权 \(L^p\) 意义下成立，其中
\[
G(y):=\frac{1}{\pi(1+y^2)},
\qquad
U_k(-y)=(-1)^kU_k(y).
\]
若 \(f\) 在 \(1\) 的邻域内属于 \(C^{2R-2}\) 且 \(f(1)=0\)，则对应的 Csisz\'ar 散度
\[
D_f(\tilde h_t\mid \tilde g_t)
:=
\int_{\RR}\tilde g_t(x)\,
f\!\left(\frac{\tilde h_t(x)}{\tilde g_t(x)}\right)\dd x
\]
具有只含偶次幂的喷流展开
\[
D_f(\tilde h_t\mid \tilde g_t)
=
\sum_{j=2}^{R}C_{2j}(f,\tilde\nu)\,t^{-2j}
+o(t^{-2R}).
\]
特别地，一切奇次阶系数 \(t^{-(2j+1)}\) 恒等为零。
\leanverified{paper\_conclusion\_poisson\_cauchy\_fdiv\_jet\_odd\_order\_vanishing}
\end{theorem}
\begin{proof}
把 \(f\) 在 \(1\) 附近作 Taylor 展开：
\[
f(1+\delta)
=
\sum_{n=1}^{2R-2}\frac{f^{(n)}(1)}{n!}\delta^n+o(\delta^{2R-2}).
\]
代入 \(\delta_t\) 的展开，并用
\[
\int_{\RR}\tilde g_t\,\delta_t\,\dd x=0
\]
消去一次项，则 \(t^{-m}\) 的系数由有限个积分项线性组合而成，这些积分项都形如
\[
\int_{\RR}G(y)\,U_{k_1}(y)\cdots U_{k_\ell}(y)\,\dd y,
\qquad
k_1+\cdots+k_\ell=m.
\]
由于 \(G\) 为偶函数，而
\[
U_{k_1}\cdots U_{k_\ell}
\]
的奇偶性恰为 \((-1)^m\)，故当 \(m\) 为奇数时，被积函数整体为奇函数，积分必为零。于是全部奇次幂系数恒等消失，只剩偶次幂。证毕。
\end{proof}

\begin{corollary}[镜像测度的完整喷流盲性]\label{cor:conclusion-poisson-cauchy-mirror-jet-blindness}
\leanverified{paper\_conclusion\_poisson\_cauchy\_mirror\_jet\_blindness}
定义关于均值 \(\bar\gamma\) 的镜像测度
\[
\tilde\nu^\vee(A):=\tilde\nu(2\bar\gamma-A).
\]
则对每个 \(k\ge 2\)，有
\[
\mu_k(\tilde\nu^\vee)=(-1)^k\mu_k(\tilde\nu).
\]
在定理 \ref{thm:conclusion-poisson-cauchy-fdiv-jet-odd-order-vanishing} 的全部适用阶内，对任意 \(C^{2R-2}\) 的 \(f\)-散度，\(\tilde\nu\) 与 \(\tilde\nu^\vee\) 的全部偶次喷流系数完全相同。并且三阶去方差总变差常数满足
\[
\lim_{t\to\infty} t^3\|\Delta_t^\vee\|_{L^1}
=
\lim_{t\to\infty} t^3\|\Delta_t\|_{L^1}
=
\frac{|\mu_3|}{\pi}.
\]
因此，整个低阶 Poisson--Cauchy 信息喷流对“绕均值镜像”完全失明；它至多能看见奇矩的绝对值，而不能看见其符号。
\end{corollary}
\begin{proof}
镜像变换直接给出
\[
\gamma-\bar\gamma\longmapsto -(\gamma-\bar\gamma),
\]
故
\[
\mu_k(\tilde\nu^\vee)=(-1)^k\mu_k(\tilde\nu).
\]
而由定理 \ref{thm:conclusion-poisson-cauchy-fdiv-jet-odd-order-vanishing}，任一 \(f\)-散度喷流的 \(t^{-2j}\) 系数只由总指标为偶数的矩单项式组成，因此在上述符号翻转下保持不变。最后由定理 \ref{thm:conclusion-poisson-cauchy-third-order-tv-constant}，
\[
\lim_{t\to\infty}t^3\|\Delta_t\|_{L^1}=\frac{|\mu_3|}{\pi},
\]
右端仅依赖于 \(|\mu_3|\)，故镜像下亦保持不变。证毕。
\end{proof}

\begin{theorem}[KL 六阶校正项的严格负性与单侧逼近律]\label{thm:conclusion-poisson-cauchy-kl-sixth-negative-one-sided-approach}
\leanverified{paper\_conclusion\_poisson\_cauchy\_kl\_sixth\_negative\_one\_sided\_approach}
设 \(\mu_2>0\)。令
\[
K_4:=\frac{\mu_2^2}{8},
\qquad
K_6:=\frac{\mu_2^3}{64}-\frac{\mu_2\mu_4}{8}+\frac{3\mu_3^2}{32}.
\]
则
\[
D_{\mathrm{KL}}(\tilde h_t\mid \tilde g_t)=K_4\,t^{-4}+K_6\,t^{-6}+o(t^{-6}),
\]
并且有严格估计
\[
K_6\le -\,\frac{7\mu_2^3+2\mu_3^2}{64}<0.
\]
从而
\[
t^4D_{\mathrm{KL}}(\tilde h_t\mid \tilde g_t)
=
K_4+K_6\,t^{-2}+o(t^{-2})
\]
从下方逼近极限 \(K_4=\mu_2^2/8\)。
\end{theorem}
\begin{proof}
六阶展开由定理 \ref{thm:conclusion-poisson-cauchy-kl-sixth-order-expansion} 给出。再由 Pearson 不等式
\[
\mu_2\mu_4\ge \mu_3^2+\mu_2^3
\]
可得
\[
K_6
=
\frac{\mu_2^3}{64}-\frac{\mu_2\mu_4}{8}+\frac{3\mu_3^2}{32}
\le
\frac{\mu_2^3}{64}-\frac{\mu_3^2+\mu_2^3}{8}+\frac{3\mu_3^2}{32}
=
-\frac{7\mu_2^3+2\mu_3^2}{64}.
\]
由于 \(\mu_2>0\)，右端严格为负，故 \(K_6<0\)。于是 \(t^4D_{\mathrm{KL}}\) 对极限 \(K_4\) 的逼近必从下方发生。证毕。
\end{proof}

\begin{theorem}[R\'enyi 六阶系数的严格负性与 \(\alpha\)-单调性]\label{thm:conclusion-poisson-cauchy-renyi-sixth-negative-monotone}
\leanverified{paper\_conclusion\_poisson\_cauchy\_renyi\_sixth\_negative\_monotone}
对任意 \(\alpha>0,\ \alpha\neq 1\)，定义六阶系数
\[
R_6(\alpha):=
\lim_{t\to\infty}
t^6\left(
D_\alpha(\tilde h_t\mid \tilde g_t)-\frac{\alpha\mu_2^2}{8t^4}
\right).
\]
则有显式公式
\[
R_6(\alpha)
=
\alpha\left(
\frac{3\mu_3^2}{32}-\frac{\mu_2\mu_4}{8}
-\frac{(\alpha-2)\mu_2^3}{64}
\right).
\]
并且：
\begin{enumerate}
\item 对一切 \(\alpha>0\)，都有 \(R_6(\alpha)<0\)；
\item 函数 \(R_6(\alpha)\) 在 \((0,\infty)\) 上严格递减；
\item 归一化系数 \(R_6(\alpha)/\alpha\) 关于 \(\alpha\) 为仿射函数，其斜率恒等于 \(-\mu_2^3/64\)。
\end{enumerate}
\end{theorem}
\begin{proof}
显式公式由定理 \ref{thm:conclusion-poisson-cauchy-renyi-sixth-order-expansion} 直接给出。仍由 Pearson 不等式，
\[
\mu_2\mu_4\ge \mu_3^2+\mu_2^3,
\]
故
\[
\frac{R_6(\alpha)}{\alpha}
\le
\frac{3\mu_3^2}{32}-\frac{\mu_3^2+\mu_2^3}{8}-\frac{(\alpha-2)\mu_2^3}{64}
=
-\frac{2\mu_3^2+(\alpha+6)\mu_2^3}{64}<0,
\]
于是 \(R_6(\alpha)<0\)。

再求导，
\[
R_6'(\alpha)
=
\frac{3\mu_3^2}{32}-\frac{\mu_2\mu_4}{8}+\frac{(1-\alpha)\mu_2^3}{32}.
\]
同样由 Pearson 不等式，
\[
R_6'(\alpha)
\le
\frac{3\mu_3^2}{32}-\frac{\mu_3^2+\mu_2^3}{8}+\frac{(1-\alpha)\mu_2^3}{32}
=
-\frac{2\mu_3^2+(2\alpha+6)\mu_2^3}{64}<0.
\]
故 \(R_6(\alpha)\) 严格递减。

最后，
\[
\frac{R_6(\alpha)}{\alpha}
=
\frac{3\mu_3^2}{32}-\frac{\mu_2\mu_4}{8}
+\frac{2\mu_2^3}{64}
-\frac{\alpha\mu_2^3}{64},
\]
显然关于 \(\alpha\) 仿射，斜率正是 \(-\mu_2^3/64\)。证毕。
\end{proof}

\begin{theorem}[\((T_3,K_4,K_6)\) 的三常数精确层析]\label{thm:conclusion-poisson-cauchy-three-constant-moment-tomography}
\leanverified{paper\_conclusion\_poisson\_cauchy\_three\_constant\_moment\_tomography}
定义三个可观测渐近常数
\[
T_3:=\lim_{t\to\infty}t^3\|\Delta_t\|_{L^1},
\qquad
K_4:=\lim_{t\to\infty}t^4D_{\mathrm{KL}}(\tilde h_t\mid \tilde g_t),
\]
\[
K_6:=\lim_{t\to\infty}t^6\left(D_{\mathrm{KL}}(\tilde h_t\mid \tilde g_t)-K_4t^{-4}\right).
\]
则它们唯一恢复
\[
\mu_2,\qquad |\mu_3|,\qquad \mu_4
\]
这三个低阶矩量，且有显式公式
\[
\mu_2=\sqrt{8K_4},
\qquad
|\mu_3|=\pi T_3,
\]
\[
\mu_4
=
\frac{3\pi^2T_3^2}{4\mu_2}
+\frac{\mu_2^2}{8}
-\frac{8K_6}{\mu_2}.
\]
因此，三个标量常数已经足以精确反演 \((\mu_2,|\mu_3|,\mu_4)\)。
\end{theorem}
\begin{proof}
由定理 \ref{thm:conclusion-poisson-cauchy-third-order-tv-constant}，
\[
T_3=\frac{|\mu_3|}{\pi}.
\]
由定理 \ref{thm:conclusion-poisson-cauchy-kl-sixth-order-expansion}，
\[
K_4=\frac{\mu_2^2}{8},
\qquad
K_6=\frac{\mu_2^3}{64}-\frac{\mu_2\mu_4}{8}+\frac{3\mu_3^2}{32}.
\]
前两式直接给出
\[
\mu_2=\sqrt{8K_4},
\qquad
|\mu_3|=\pi T_3.
\]
把 \(\mu_3^2=\pi^2T_3^2\) 与 \(\mu_2=\sqrt{8K_4}\) 代回 \(K_6\) 的表达式并解出 \(\mu_4\)，即得所示公式。证毕。
\end{proof}

\begin{theorem}[\(C^3\) Csisz\'ar 散度六阶系数的 jet--仿射律]\label{thm:conclusion-poisson-cauchy-c3-fdiv-sixth-jet-affine-law}
设 \(f\) 在 \(1\) 的邻域内属于 \(C^3\)，且 \(f(1)=0\)。则其对应的 Csisz\'ar 散度满足
\[
D_f(\tilde h_t\mid \tilde g_t)
=
\frac{f''(1)\mu_2^2}{8}\,t^{-4}
+C_6(f)\,t^{-6}
+o(t^{-6}),
\]
其中
\[
C_6(f)
=
f''(1)\left(\frac{3\mu_3^2}{32}-\frac{\mu_2\mu_4}{8}\right)
-\frac{f^{(3)}(1)\mu_2^3}{64}.
\]
特别地，六阶系数只依赖于 \(f\) 在 \(1\) 处的二、三阶 jet
\[
\bigl(f''(1),f^{(3)}(1)\bigr).
\]
\end{theorem}
\begin{proof}
对 \(f(1+\delta_t)\) 在 \(\delta_t=0\) 处作三阶 Taylor 展开：
\[
f(1+\delta_t)
=
f'(1)\delta_t+\frac{f''(1)}{2}\delta_t^2+\frac{f^{(3)}(1)}{6}\delta_t^3+O(\delta_t^4).
\]
质量守恒给出
\[
\int_{\RR}\tilde g_t\,\delta_t\,\dd x=0,
\]
且 \(\delta_t=O(t^{-2})\)，故
\[
D_f(\tilde h_t\mid \tilde g_t)
=
\frac{f''(1)}{2}\int \tilde g_t\delta_t^2\,\dd x
+\frac{f^{(3)}(1)}{6}\int \tilde g_t\delta_t^3\,\dd x
+O(t^{-8}).
\]
再用宿主中 KL 六阶展开证明所列积分常数
\[
\int G\,U_2^2\,\dd y=\frac14,\qquad
\int G\,U_3^2\,\dd y=\frac{3}{16},\qquad
\int G\,U_2U_4\,\dd y=-\frac18,\qquad
\int G\,U_2^3\,\dd y=-\frac{3}{32},
\]
以及
\[
\delta_t=\frac{\mu_2}{t^2}U_2+\frac{\mu_3}{t^3}U_3+\frac{\mu_4}{t^4}U_4+O(t^{-5}),
\]
便得
\[
\int \tilde g_t\delta_t^2\,\dd x
=
\frac{\mu_2^2}{4}t^{-4}
+\left(\frac{3\mu_3^2}{16}-\frac{\mu_2\mu_4}{4}\right)t^{-6}
+o(t^{-6}),
\]
\[
\int \tilde g_t\delta_t^3\,\dd x
=
-\frac{3\mu_2^3}{32}t^{-6}+o(t^{-6}).
\]
代回并整理即得所述公式。证毕。
\end{proof}

\begin{theorem}[六阶 universality class 的双散度判别]\label{thm:conclusion-poisson-cauchy-sixth-universality-class}
\leanverified{paper\_conclusion\_poisson\_cauchy\_sixth\_universality\_class}
设 \(\tilde\nu,\tilde\nu'\) 为两条中心化扰动测度。则以下两命题等价：
\begin{enumerate}
\item 它们满足
\[
\mu_2(\tilde\nu)=\mu_2(\tilde\nu'),\qquad
|\mu_3(\tilde\nu)|=|\mu_3(\tilde\nu')|,\qquad
\mu_4(\tilde\nu)=\mu_4(\tilde\nu');
\]
\item 它们具有相同的三阶去方差总变差常数，且对任意两个 \(C^3\) 散度生成元 \(f,g\) 满足
\[
\det
\begin{pmatrix}
f''(1) & f^{(3)}(1)\\
g''(1) & g^{(3)}(1)
\end{pmatrix}\neq0,
\]
相应的 \(t^{-4}\) 与 \(t^{-6}\) 系数全部一致。
\end{enumerate}
亦即，\((\mu_2,|\mu_3|,\mu_4)\) 正是六阶信息喷流的完整 universality coordinates。
\end{theorem}
\begin{proof}
\((1)\Rightarrow(2)\) 显然：由定理 \ref{thm:conclusion-poisson-cauchy-c3-fdiv-sixth-jet-affine-law}，任意 \(C^3\) 散度的 \(t^{-4}\) 与 \(t^{-6}\) 系数只依赖于 \(\mu_2,\mu_3^2,\mu_4\)；而三阶总变差常数只依赖于 \(|\mu_3|\)。

\((2)\Rightarrow(1)\)。先由任一散度的 \(t^{-4}\) 系数相等得
\[
\mu_2(\tilde\nu)=\mu_2(\tilde\nu').
\]
再由三阶总变差常数相等和定理 \ref{thm:conclusion-poisson-cauchy-third-order-tv-constant} 得
\[
|\mu_3(\tilde\nu)|=|\mu_3(\tilde\nu')|.
\]
最后，把两条生成元的六阶系数写成线性系统：
\[
\begin{pmatrix}
C_6(f)\\[1mm]
C_6(g)
\end{pmatrix}
=
\begin{pmatrix}
f''(1) & -f^{(3)}(1)/64\\
g''(1) & -g^{(3)}(1)/64
\end{pmatrix}
\begin{pmatrix}
\frac{3\mu_3^2}{32}-\frac{\mu_2\mu_4}{8}\\[1mm]
\mu_2^3
\end{pmatrix}.
\]
系数矩阵可逆，故若两条测度对应的六阶系数相同，则
\[
\frac{3\mu_3^2}{32}-\frac{\mu_2\mu_4}{8}
\]
也相同。前面已知 \(\mu_2\) 与 \(\mu_3^2\) 相同，因此 \(\mu_4\) 相同。证毕。
\end{proof}

\begin{theorem}[八阶 \(f\)-jet 的三维线性层析]\label{thm:conclusion-poisson-cauchy-eighth-jet-linear-tomography}
在八阶喷流的适用设定下，定义三维 moment--jet 向量
\[
\mathcal J(\tilde\nu):=
\begin{pmatrix}
24\mu_2\mu_6-60\mu_3\mu_5+40\mu_4^2\\[1mm]
24\mu_2^2\mu_4-18\mu_2\mu_3^2\\[1mm]
3\mu_2^4
\end{pmatrix}.
\]
则任意 \(C^4\) 的 Csisz\'ar 散度生成元 \(f\) 的八阶系数满足严格线性律
\[
512\,C_8(f)
=
\begin{pmatrix}
f''(1)& f^{(3)}(1)& f^{(4)}(1)
\end{pmatrix}
\mathcal J(\tilde\nu).
\]
特别地，若取三个 \(C^4\) 生成元 \(f_1,f_2,f_3\) 使
\[
\det
\begin{pmatrix}
f_1''(1)&f_1^{(3)}(1)&f_1^{(4)}(1)\\
f_2''(1)&f_2^{(3)}(1)&f_2^{(4)}(1)\\
f_3''(1)&f_3^{(3)}(1)&f_3^{(4)}(1)
\end{pmatrix}\neq0,
\]
则三条八阶系数
\[
C_8(f_1),\qquad C_8(f_2),\qquad C_8(f_3)
\]
唯一恢复整个向量 \(\mathcal J(\tilde\nu)\)。
\end{theorem}
\begin{proof}
这正是定理 \ref{thm:conclusion-poisson-cauchy-fdiv-eighth-order-jet} 的八阶闭式改写。该定理给出
\[
C_8(f)
=
\frac1{512}\Bigl(
f''(1)\bigl(24\mu_2\mu_6-60\mu_3\mu_5+40\mu_4^2\bigr)
+f^{(3)}(1)\bigl(24\mu_2^2\mu_4-18\mu_2\mu_3^2\bigr)
+f^{(4)}(1)\cdot 3\mu_2^4
\Bigr).
\]
把右端视作 \((f''(1),f^{(3)}(1),f^{(4)}(1))\) 对 \(\mathcal J(\tilde\nu)\) 的线性配对，即得第一式。

若三条生成元的 jet 矩阵可逆，则由三条线性方程
\[
512
\begin{pmatrix}
C_8(f_1)\\
C_8(f_2)\\
C_8(f_3)
\end{pmatrix}
=
\begin{pmatrix}
f_1''(1)&f_1^{(3)}(1)&f_1^{(4)}(1)\\
f_2''(1)&f_2^{(3)}(1)&f_2^{(4)}(1)\\
f_3''(1)&f_3^{(3)}(1)&f_3^{(4)}(1)
\end{pmatrix}
\mathcal J(\tilde\nu)
\]
即可唯一恢复 \(\mathcal J(\tilde\nu)\)。证毕。
\end{proof}

\begin{proposition}[固定一阶矩类中的二阶可见模空间是一维]\label{prop:conclusion-poisson-second-order-visible-moduli-one-dimensional}
\leanverified{paper\_conclusion\_poisson\_second\_order\_visible\_moduli\_one\_dimensional}
设 \(G_t,\widetilde G_t\) 为两条点状缺陷测度的 Poisson 粗粒化剖面，并假定它们具有相同总质量 \(M_0>0\) 与相同一阶矩 \(M_1\)。记共同重心为
\[
\bar\gamma:=\frac{M_1}{M_0},
\]
并令 \(\tilde\nu,\tilde\nu'\) 为相应的归一化权重测度，其方差分别记为
\[
\mathrm{Var}:=\int (\gamma-\bar\gamma)^2\,\dd\tilde\nu(\gamma),
\qquad
\widetilde{\mathrm{Var}}:=\int (\gamma-\bar\gamma)^2\,\dd\tilde\nu'(\gamma).
\]
则存在一致展开
\[
G_t(x)-\widetilde G_t(x)
=
\frac{M_0\bigl(\mathrm{Var}-\widetilde{\mathrm{Var}}\bigr)}{\pi}\,
\frac{t\bigl(3(x-\bar\gamma)^2-t^2\bigr)}{\bigl((x-\bar\gamma)^2+t^2\bigr)^3}
+O(t^{-4}),
\]
其中余项在 \(x\in\RR\) 上一致。因而，在固定 \((M_0,M_1)\) 的等矩类内，Poisson 粗粒化在二阶精度上的全部可见差异，仅由单个标量 \(\mathrm{Var}\) 参数化。
\end{proposition}
\begin{proof}
由命题 \ref{con:xi-poisson-coarsegraining-second-order-normal-form}，对每一条归一化权重测度都有展开
\[
h_t(x)=g_t(x)+\mathrm{Var}\cdot r_t(x)+O(t^{-4}),
\qquad
\widetilde h_t(x)=g_t(x)+\widetilde{\mathrm{Var}}\cdot r_t(x)+O(t^{-4}),
\]
其中
\[
g_t(x):=\frac{1}{\pi}\frac{t}{(x-\bar\gamma)^2+t^2},
\qquad
r_t(x):=\frac{1}{\pi}\frac{t\bigl(3(x-\bar\gamma)^2-t^2\bigr)}{\bigl((x-\bar\gamma)^2+t^2\bigr)^3}.
\]
而原始粗粒化剖面满足
\[
G_t=M_0h_t,
\qquad
\widetilde G_t=M_0\widetilde h_t.
\]
两式相减即得
\[
G_t(x)-\widetilde G_t(x)
=
M_0\bigl(\mathrm{Var}-\widetilde{\mathrm{Var}}\bigr)r_t(x)+O(t^{-4}),
\]
这正是所需展开。由右端可见，在固定 \((M_0,M_1)\) 的条件下，二阶可见差异只通过 \(\mathrm{Var}-\widetilde{\mathrm{Var}}\) 进入，故可见模空间是一维的。证毕。
\end{proof}

\begin{corollary}[二阶剖面的双结点盲区与单点反演]\label{cor:conclusion-poisson-second-order-two-node-blindness}
\leanverified{paper\_conclusion\_poisson\_second\_order\_two\_node\_blindness}
在命题 \ref{prop:conclusion-poisson-second-order-visible-moduli-one-dimensional} 的单测度情形下，令
\[
g_t(x):=\frac{M_0}{\pi}\frac{t}{(x-\bar\gamma)^2+t^2},
\qquad
y:=\frac{x-\bar\gamma}{t},
\]
并定义二阶归一化剖面
\[
\mathcal R_t(y):=
t^2\left(\frac{G_t(\bar\gamma+ty)}{g_t(\bar\gamma+ty)}-1\right).
\]
则有一致极限
\[
\mathcal R_t(y)=\frac{\mathrm{Var}}{2}\,\mathcal H(y)+O(t^{-1}),
\qquad
\mathcal H(y):=\frac{2(3y^2-1)}{(1+y^2)^2}.
\]
因此 \(\mathcal H\) 恰有两个零点
\[
y=\pm \frac{1}{\sqrt3}.
\]
特别地：
\begin{enumerate}
\item 二阶方差信号只有这两个相似变量位置完全不可见；
\item 对任意 \(y_0\neq \pm 1/\sqrt3\)，单点观测已可恢复方差，
\[
\mathrm{Var}
=
\lim_{t\to\infty}\frac{2\mathcal R_t(y_0)}{\mathcal H(y_0)}.
\]
\end{enumerate}
换言之，Poisson--Cauchy 二阶校正的观测盲区并非连续族，而是严格退化为两个离散结点。
\end{corollary}
\begin{proof}
由推论 \ref{cor:xi-cauchy-poisson-density-ratio-second-order-profile}，
\[
\mathcal R_t(y)=\frac{\mathrm{Var}}{2}\,\mathcal H(y)+O(t^{-1})
\qquad\text{一致于 }y\in\RR.
\]
而定理 \ref{thm:xi-cauchy-poisson-second-order-shape-limit-node-rigidity} 已给出 \(\mathcal H(y)\) 的零点仅为
\[
y=\pm \frac{1}{\sqrt3}.
\]
故除这两个结点外，\(\mathcal H(y_0)\neq 0\)，从而
\[
\mathrm{Var}
=
\lim_{t\to\infty}\frac{2\mathcal R_t(y_0)}{\mathcal H(y_0)}.
\]
这便得到二阶盲区与单点反演公式。证毕。
\end{proof}

\noindent
以上结果把 Poisson--Cauchy 信息渐近压成一条由低阶至高阶的喷流几何链：在二阶层面，固定 \((M_0,M_1)\) 的粗粒化可见模空间严格退化为单一方差参数，并且观测盲区只剩 \(\pm 1/\sqrt3\) 两个离散结点；在更高阶层面，轮廓函数的奇偶性又强制一切 \(f\)-散度的奇次喷流全局消隐，并由此诱导出绕均值镜像的完整喷流盲性。进一步，KL 与 R\'enyi 的六阶校正项都被迫为严格负项，从而六阶主导常数总以单侧方式逼近其四阶极限；\((\mu_2,|\mu_3|,\mu_4)\) 则作为三常数层析坐标完全控制全部六阶 universality class；最后，所有八阶 \(f\)-散度又统一坍缩为三维 moment--jet 向量的线性层析问题。于是，\S\ref{subsec:conclusion-poisson-cauchy-high-order-entropy-fiber-rigidity} 中分散的二阶与高阶展开在结论层被压接为同一条 parity--jet--tomography 刚性链。

\endinput
